% Vista científica exacta materializada desde una rama condicional.
% Fuente: source/demostraciones_primarias/31_06c_sectores_topologicos_no_omision.tex; rama: HMTIncludeCMBlock.
% No modifica el contenido matemático de la rama de origen.
\chapter[Estructura compleja \texorpdfstring{\(\widetilde J_M^{\,2}=-I\)}{J M al cuadrado = -I} y sectores topológicos]
{Estructura compleja \texorpdfstring{\(\widetilde J_M^{\,2}=-I\)}{J M al cuadrado = -I} sobre el estrato fijo, involuciones dodecafásicas y sectores topológicos}
\label{chap:sectores-topologicos-no-omision}
\index{sector topológico}\index{Fibonacci, anillo basado}
\index{Majorana}\index{Ising}\index{grupo de trenzas}

La involución \(C_M\), el anillo basado asociado a \(F_{\rm av}\), los
caracteres de trenza y las categorías apuntadas pertenecen a estratos
algebraicos distintos. Se presentan con sus dominios y controles de
identificación para que una coincidencia nominal no sustituya al morfismo que
debería relacionarlos. El censo del atlas, la jerarquía espectral infinita
cuyas ocho alturas históricas son testigos iniciales, y las completaciones
estadísticas conservan las demostraciones de
\cref{thm:censo-atlas-81,thm:descenso-necesario-fav,%
thm:principio-exclusion-pauli,thm:dos-cuatro-pi-estadistica}.

Cada resultado declara directamente su dominio, su construcción y su
conclusión; ninguna coincidencia nominal modifica esas dependencias.

\section{Estrato periódico de la involución dodecafásica}
\label{sec:involucion-cm-censo}

El dominio \(\mathscr T_{12}^{\rm word}\), la conjugación
\(C_M\mathbf t=-\operatorname{rev}(\mathbf t)\), su carácter involutivo y el
censo \(\lvert\operatorname{Fix}(C_M)\rvert=3^6=729\) tienen su residencia
primaria en las
\cref{def:palabras-cm,thm:involucion-cm-censo}, dentro de la banda de
partícula. Aquí se añade únicamente la condición periódica necesaria para el
estrato topológico. Si \(r\) denota el desplazamiento cíclico
\((r\mathbf t)_m=t_{m-1}\), ponemos
\[
 \operatorname{Per}_6
 :=\{\mathbf t\in\mathscr T_{12}^{\rm word}:r^6\mathbf t=\mathbf t\}.
\]

\begin{proposition}[Censo del estrato periódico]
\label{prop:estrato-periodico-cm-rev7}
Al imponer periodicidad de período divisor de seis,
\[
 \bigl|\operatorname{Fix}(C_M)\cap\operatorname{Per}_6\bigr|
 =3^3=27.
\]
Esta identidad es exacta.
\end{proposition}

\begin{proof}
El \cref{thm:involucion-cm-censo} da la forma general de los puntos fijos.
Si además \(r^6\mathbf t=\mathbf t\), la segunda mitad de la palabra coincide
con la primera. Escribiendo las seis primeras coordenadas, las dos
condiciones equivalen a
\[
 (t_0,t_1,t_2,t_3,t_4,t_5)
 =(a,b,c,-c,-b,-a).
\]
Los parámetros \(a,b,c\) son independientes y toman tres valores; por tanto,
la intersección tiene \(3^3\) elementos.
\end{proof}

El ejemplo general de un punto fijo, antes de exigir período seis, es
\[
 (a,b,c,d,e,f,-f,-e,-d,-c,-b,-a).
\]
Esta forma coordinada permite comprobar a mano tanto la involución como el
conteo y evita interpretar el entero \(729\) mediante otro objeto homónimo.

\begin{remark}[Censo exhaustivo de la involución]
\label{rem:procedencia-n40-cm}
La enumeración de las \(3^{12}=531\,441\) palabras no encuentra ninguna
violación de la involución y obtiene
\[
 |\operatorname{Fix}(C_M)|=729,
 \qquad
 |\operatorname{Fix}(C_M)\cap\operatorname{Per}_6|=27.
\]
La demostración primaria y la \cref{prop:estrato-periodico-cm-rev7} explican
ambos números sin depender del recorrido: los puntos fijos tienen seis
coordenadas libres, y la condición adicional de período seis reduce la
elección a tres. El censo exhaustivo confirma que no
existen órbitas excepcionales fuera de esa descripción.
\end{remark}

\begin{criterion}[Falsadores del censo \(C_M\)]
\label{crit:falsadores-cm}
Refutaría el resultado cualquiera de los hechos siguientes: una palabra para
la cual \(C_M^2\mathbf t\ne\mathbf t\); un punto fijo que no posea la forma
coordenada anterior; un total exhaustivo distinto de \(729\); o un total
distinto de \(27\) después de imponer \(r^6\mathbf t=\mathbf t\). La
coincidencia nominal de este \(729\) con otro cardinal no identifica sus
dominios.
\end{criterion}

\section{Estructura cuadrática transportada al estrato fijo}
\label{sec:estructura-cuadratica-estrato-fijo-cm}

Identificamos el alfabeto de signos con el cuerpo ternario mediante
\[
 \{-1,0,+1\}\xrightarrow{\sim}\mathbb F_3,
 \qquad
 -1\longmapsto2,\quad 0\longmapsto0,\quad +1\longmapsto1.
\]
Con esta identificación,
\(\mathscr T_{12}^{\rm word}\cong\mathbb F_3^{12}\), y tanto
\(C_M=-\operatorname{rev}\) como el desplazamiento \(r\) son operadores
\(\mathbb F_3\)-lineales. En consecuencia,
\(\operatorname{Fix}(C_M)\) y
\(\operatorname{Per}_6=\ker(r^6-I)\) son subespacios vectoriales.

La forma coordinada de los puntos fijos define el isomorfismo lineal
\[
 s_{\rm fix}:\mathbb F_3^6\xrightarrow{\sim}\operatorname{Fix}(C_M),
\]
\[
 s_{\rm fix}(a,b,c,d,e,f)
 =(a,b,c,d,e,f,-f,-e,-d,-c,-b,-a).
\]
Sea \(J_W\) el operador de Paley--Witt construido en el
\cref{thm:f9-tres-desde-paley}, con la base, la polarización y la orientación
allí declaradas.

\begin{proposition}[Estructura compleja \(\widetilde J_M^{\,2}=-I\) sobre el estrato fijo]
\label{prop:raiz-menos-uno-estrato-fijo-cm}
El transporte
\[
\widetilde J_M:=s_{\rm fix}J_Ws_{\rm fix}^{-1}
\in\operatorname{End}_{\mathbb F_3}
   \bigl(\operatorname{Fix}(C_M)\bigr)
\]
satisface
\[
 \widetilde J_M^{\,2}=-I.
\]
\end{proposition}

\begin{proof}
Por conjugación,
\[
 \widetilde J_M^{\,2}
 =s_{\rm fix}J_Ws_{\rm fix}^{-1}
  s_{\rm fix}J_Ws_{\rm fix}^{-1}
 =s_{\rm fix}J_W^2s_{\rm fix}^{-1}
 =-I.
\]
\end{proof}

El resultado está demostrado con la carta \(s_{\rm fix}\), la base, la polarización y
la orientación como estructura de partida explícita. No identifica
\(C_M\) con \(J_W\): el primero es una involución del espacio de palabras y
el segundo una estructura cuadrática sobre su estrato fijo. Tampoco demuestra
que el transporte sea natural respecto de toda la dinámica TPK.

Pongamos
\[
 W_{27}
 :=\operatorname{Fix}(C_M)\cap\operatorname{Per}_6.
\]
La inclusión
\[
 W_{27}\hookrightarrow\operatorname{Fix}(C_M)
\]
es canónica, y la \cref{prop:estrato-periodico-cm-rev7} prueba que
\(\dim_{\mathbb F_3}W_{27}=3\).

\begin{corollary}[Obstrucción cuadrática en el estrato de veintisiete]
\label{cor:obstruccion-cuadratica-w27}
El subespacio \(W_{27}\) no puede ser estable bajo
\(\widetilde J_M\).
\end{corollary}

\begin{proof}
Si fuera estable, la relación
\(\widetilde J_M^2=-I\) lo convertiría en un módulo sobre
\[
 \mathbb F_3[X]/(X^2+1)\cong\mathbb F_9.
\]
Todo espacio vectorial sobre \(\mathbb F_9\) tiene dimensión par cuando se
considera sobre \(\mathbb F_3\), en contradicción con
\(\dim_{\mathbb F_3}W_{27}=3\).
\end{proof}

La igualdad cardinal entre \(W_{27}\) y las veintisiete rectas de Schläfli
no proporciona un mapa entre ambos objetos. Tal relación exigiría construir
una aplicación de incidencia que preserve la orientación central y la acción
del estabilizador, además de una clausura dinámica compatible. No se
interpone el atlas de \(81\) celdas entre \(W_{27}\) y el estrato fijo de
\(729\) estados: el atlas de partículas pertenece a otro dominio.

\section[Niveles de Majorana]{Estratificación matemática del objeto Majorana en cuatro niveles}
\label{sec:cuatro-niveles-majorana}

La inversión \(C_M\), la paridad exterior, un fermión relativista de
Majorana y un modo cero de Majorana son objetos de tipos distintos:

\begin{enumerate}
 \item \(C_M=-\operatorname{rev}\) actúa sobre
 \(\mathscr T_{12}^{\rm word}\). Su involutividad y su censo se demuestran
 exactamente; esta construcción no proporciona por sí misma un espinor, una
 ecuación de campo ni un término de masa.

 \item El carácter de signo \(q=-1\) y la completación exterior pertenecen
 al álgebra de intercambio. Relativamente a CAR y a la graduación declarada,
 producen \(N_s^2=N_s\), espectro \(\{0,1\}\) y el retorno de orientación
 \(2\pi/4\pi\). Estas identidades se deducen de CAR y de la graduación; no se
 deducen del
 mero censo de \(C_M\).

 \item Un fermión relativista de Majorana es un objeto espinorial provisto de
 representación relativista, conjugación de carga compatible, condición de
 autorconjugación, término de masa y dinámica. Esos datos no forman parte de
 la definición de \(C_M\) ni se deducen de su censo.

 \item Un modo cero de Majorana, o equivalentemente el sector no abeliano de
 una teoría de Ising en la acepción topológica, requiere simples
 \(\mathbf 1,\psi,\sigma\) con
 \[
  \psi\otimes\psi\simeq\mathbf 1,\qquad
  \psi\otimes\sigma\simeq\sigma,\qquad
  \sigma\otimes\sigma\simeq\mathbf 1\oplus\psi,\qquad
  d_\sigma=\sqrt2,
 \]
 junto con asociador y trenzado coherentes. La brecha, el desdoblamiento y la
 plataforma pertenecen a una realización física adicional; las reglas de
 fusión anteriores no definen una masa de polo universal.
\end{enumerate}

La separación es de dominio y no de nomenclatura: involución de palabras,
paridad CAR, espinor autorconjugado y simple no abeliano son cuatro objetos
matemáticos distintos. En ausencia de un morfismo explícito entre ellos, no
se identifica ninguno por compartir el nombre «Majorana».
